\section{Future work}
\label{sec:future-work}

The system runs end to end on real hardware: record, label, train, deliver over the air, classify
on-device, and display live. The following items are known and planned, roughly in priority order.

\begin{description}[leftmargin=1.6em,style=nextline]
  \item[Explicit ``neither'' class] A still or resting cat is already not forced into a behaviour: the
        rules-based activity gate logs it as low-power rest, and short or uncertain windows report
        unknown. The remaining refinement is to train an explicit ``neither'' class so an \emph{awake}
        cat doing something untrained is not pushed into the nearest trained behaviour , a data and
        modelling task, not a firmware change (\cref{sec:training-pipeline}).
  \item[Power tuning] Tune the rules-based activity gate's motion thresholds and rest hold-off, and
        the audio path's duty cycle, against the measured power budget (\cref{sec:evaluation}) to
        extend collar runtime. A deeper step is true hardware wake-on-motion (an IMU interrupt waking
        the SoC from System OFF) for the lowest rest current.
  \item[More behaviours] Extend the captured and labelled set beyond eat/drink/purr to scratching,
        litter-tray use, grooming, and play, exercising the label-agnostic pipeline at scale.
  \item[Cross-cat generalisation] Multiple cats per household are already supported (each wears its own
        collar, with a dashboard switcher and per-cat data). The open problem is \emph{model}
        generalisation , training on one cat and performing well on another , since the accuracy
        evaluation is currently single-cat and within-subject (\cref{sec:evaluation}).
  \item[Firmware over-the-air] Today only models are delivered over the air; extend the mechanism to
        full firmware updates.
  \item[Publication] Write up the confidence-gated cascade as a short preprint and a venue paper.
\end{description}
